\section*{Bab \ref{ch:g1:measure} --- Mengukur dan Waktu}

\begin{solution}{exo:g1:measure:1}
Jawabannya bergantung pada pensilnya; yang paling panjang menjulur
paling jauh ketika ujungnya diimpitkan, yang paling pendek menjulur
paling dekat.
\end{solution}

\begin{solution}{exo:g1:measure:2}
Jawaban beragam; setiap panjang dibaca dengan meletakkan angka $0$
penggaris di awal benda lalu membaca sentimeter bulat di ujungnya.
\end{solution}

\begin{solution}{exo:g1:measure:3}
Sesudah Senin: Selasa. Sebelum Minggu: Sabtu. Dua hari sesudah
Jumat: Minggu.
\end{solution}

\begin{solution}{exo:g1:measure:4}
Kemarin: Rabu. Besok: Jumat. Tiga hari lagi: Minggu.
\end{solution}

\begin{solution}{exo:g1:measure:5}
Pada pukul $6$ jarum pendek menunjuk angka $6$ dan jarum panjang
menunjuk lurus ke atas ke angka $12$.
\end{solution}

\begin{solution}{exo:g1:measure:6}
Jarum pendek di $5$: pukul $5$. Jarum pendek di $11$: pukul $11$.
\end{solution}

\begin{solution}{exo:g1:measure:7}
Untuk $8$, misalnya $5 + 2 + 1$ atau $2 + 2 + 2 + 2$ (masih ada
jawaban benar yang lain).
\end{solution}

\begin{solution}{exo:g1:measure:8}
$5 + 2 + 2 + 2 = 11$. Rudi punya $11$.
\end{solution}

\begin{solution}{exo:g1:measure:9}
Dari $3$ naik ke $5$ ada $2$: Zahra menerima $2$ kembali.

\end{solution}
